\chapter[Referêncial Teórico]{Referêncial Teórico}
\label{cap_exemplos}

A revisão da literatura deve ser composta por aspectos conceituais do seu problema, para que possamos entender as definições tradicionais ou conforme a vertente teórica que você está seguindo em seu trabalho. Além disso, é importante apresentar uma retrospectiva histórica que permita compreender o cenário desse problema, especialmente em nível global, continental, nacional e local, considerando a área pesquisada.

Você também deve refletir sobre quais são os fatores que contribuem para esse problema e quais são as suas consequências. Com essa estrutura organizacional, é possível encontrar, na literatura, informações que irão elucidar os aspectos essenciais do problema, proporcionando maior clareza sobre o seu sentido prático.

\section{Aspectos Conceituais do Problema}
Nesta seção, apresentam-se as definições fundamentais que embasam o problema estudado, bem como as principais vertentes teóricas utilizadas:

\begin{itemize}
    \item \textbf{Definições Tradicionais}: Conceitos estabelecidos na literatura por autores clássicos e contemporâneos.
    \item \textbf{Vertentes Teóricas}: Diferentes abordagens e interpretações do problema, conforme a perspectiva de autores relevantes.
\end{itemize}

Esses elementos são essenciais para **clareza conceitual** e para delimitar o escopo da pesquisa.

\section{Retrospectiva Histórica}
A retrospectiva histórica é importante para compreender a evolução do problema no tempo e no espaço. Essa análise é dividida em quatro níveis:

\subsection{Cenário Global}
Aqui, discutem-se as origens e o desenvolvimento do problema em nível mundial, destacando marcos históricos, eventos e contextos globais.

\subsection{Cenário Continental}
Nesta subseção, aborda-se o problema no contexto de um continente específico, destacando suas particularidades e impactos regionais.

\subsection{Cenário Nacional}
Analisa-se o problema no contexto nacional, levando em consideração políticas públicas, dados relevantes e estudos locais.

\subsection{Cenário Local}
Por fim, examina-se o problema na localidade específica do estudo, identificando as condições locais que influenciam sua manifestação e magnitude.

\section{Fatores Contribuintes}
Os principais fatores que contribuem para o problema são discutidos a seguir:

\begin{itemize}
    \item \textbf{Fatores Sociais}: Aspectos relacionados à desigualdade, políticas públicas e organização social.
    \item \textbf{Fatores Ambientais}: Influências das condições naturais, como mudanças climáticas ou desmatamento.
    \item \textbf{Fatores Tecnológicos}: Limitações técnicas, ausência de inovação ou desafios de implementação.
\end{itemize}

Compreender esses fatores é crucial para identificar as causas e os desafios envolvidos.

\section{Consequências do Problema}
Nesta seção, apresentam-se as principais consequências práticas do problema:

\begin{itemize}
    \item \textbf{Impactos Sociais}: Efeitos diretos sobre a qualidade de vida da população.
    \item \textbf{Impactos Econômicos}: Custos financeiros e impactos na produtividade.
    \item \textbf{Impactos Científicos}: Desafios e lacunas para o avanço do conhecimento científico e tecnológico.
\end{itemize}

Essas consequências reforçam a **relevância prática** do estudo e justificam a necessidade de soluções efetivas.
